\documentclass[11pt,a4paper]{article}
\usepackage{fontspec,amsmath,amssymb,amsthm,microtype}
\setmainfont{lmroman10-regular.otf}[BoldFont=lmroman10-bold.otf,ItalicFont=lmroman10-italic.otf,BoldItalicFont=lmroman10-bolditalic.otf]
\usepackage[margin=25mm]{geometry}
\usepackage[unicode,hidelinks]{hyperref}
\hypersetup{pdfauthor={LangevinLean Project}}
\setlength{\emergencystretch}{3em}
\setlength{\parskip}{0.4em}
\setlength{\parindent}{0pt}
\allowdisplaybreaks[2]
\raggedbottom
\title{\texorpdfstring{From scalar coordinate laws to independent Gaussian vectors}{From scalar coordinate laws to independent Gaussian vectors}}
\author{LangevinLean Project}
\date{}
\begin{document}\maketitle
Draft mathematical exposition; these notes may contain mistakes.


Assume the scalar innovation array is jointly independent and each entry has law \(\mathcal N(0,1)\). Group the finitely many spatial entries at time \(n\) into a vector \(W_n\). Its coordinate product law is the standard Euclidean Gaussian law \(\mathcal N(0,I)\). Grouping disjoint subsets of an independent array preserves independence, so the vectors \(W_n\) are independent over time.



The resulting vector-sequence law is the product of these standard Gaussian laws. This conclusion uses joint independence of the scalar array, not only its coordinate marginal distributions.


\section{\texorpdfstring{The law of the scalar innovation array}{The law of the scalar innovation array}}


Regroup the jointly independent scalar coordinates indexed by \((n,i)\) into vector blocks \(W_n\). A finite block with independent standard Gaussian coordinates has the standard Euclidean Gaussian law. Product-measure identification and reindexing then show that the full vector path has law \(\bigotimes_n\mathcal N(0,I)\).



The path-level product law specifies every finite array of innovation coordinates. It therefore determines the joint law of their sequences, including dependence across time.


\section{\texorpdfstring{The standard Gaussian vector law}{The standard Gaussian vector law}}


Project the established whole-path product law to its \(n\)-th vector coordinate. Product marginals give \(P\circ W_n^{-1}=\mathcal N(0,I)\). This derives the multivariate law from the scalar coordinate assumptions rather than adding it as an algorithm hypothesis.



At a fixed index, the independent scalar standard Gaussian coordinates form the Euclidean standard Gaussian vector. The product characterization identifies the vector law.


\section{\texorpdfstring{Independence of the vector innovations}{Independence of the vector innovations}}


The coordinate maps of the vector product law are independent. Pull that independence back through the innovation-path map using the product law of the complete scalar innovation array. Thus the vector-valued \(W_n\) are independent over time, not just uncorrelated.



Disjoint time-index groups of an independent scalar array remain independent after assembly into vectors. Hence the vector innovations form an independent family with the required common Gaussian law.


\end{document}
